\chapter{Full Blade Pipeline: windIO to BeamDyn and Back}
\label{ch:windio}

A wind-turbine blade is designed as a three-dimensional laminated shell but
analysed, aeroelastically, as a one-dimensional beam. This chapter walks the
whole round trip on the IEA-22 blade: a windIO blade description becomes one
1-D shell structure gene per span station, each station is homogenized to a
Timoshenko $6\times6$, the spanwise table of those matrices is what a 1-D beam
solver needs, the beam run returns sectional force and moment resultants
station by station, and those resultants drive the dehomogenization back down
to three-dimensional stress in the wall.

Three separate tools are involved, and the boundary between them is the single
most important thing to understand before running anything. OpenSG never runs
the blade dynamics. It produces the sectional constitutive law that the beam
solver needs, and it consumes the sectional loads that the beam solver returns.

\begin{table}[htbp]
\centering
\small
\caption{The five stages of the round trip, and which tool owns each one.}
\label{tab:windio:stages}
\begin{tabular}{clll}
\toprule
Stage & What happens & Tool & Driver \\
\midrule
1 & windIO blade $\rightarrow$ per-station 1-D shell SG YAML & OpenSG\_io & \texttt{scripts/opensg\_io.py} \\
2 & 1-D shell SG $\rightarrow$ Timoshenko $6\times6$ & OpenSG-2.0 & \texttt{1\_homo\_beam\_props.py} \\
  &  &  & \texttt{4\_sweep\_51\_stations.py} \\
3 & spanwise $6\times6$ table $\rightarrow$ blade dynamics & BeamDyn / OpenFAST & run by you \\
4 & sectional loads returned as a table & BeamDyn / OpenFAST & \texttt{ff51\_rmc\_reform.dat} \\
5 & loads $+\ 6\times6$ $\rightarrow$ 3-D wall stress & OpenSG-2.0 & \texttt{2\_dehom\_v2\_vabs.py} \\
  &  &  & \texttt{4\_sweep\_51\_stations.py} \\
  & correctness check on transverse shear & OpenSG-2.0 & \texttt{3\_q\_consistency\_check.py} \\
\bottomrule
\end{tabular}
\end{table}

\section{The worked example}
\label{sec:windio:example}

Everything in this chapter runs in the directory
\texttt{examples/OpenSG\_shell/IEA\_blade\_beam/}. Run the scripts from inside
that directory: the package holds no path assumptions, and the example owns
every path it uses.

The three single-station scripts all analyse the same section,
\texttt{iea\_s10}, the IEA-22 station at $r/R = 0.20$. Its 1-D shell SG,
\texttt{iea\_s10\_shell.yaml}, is a ring of \textbf{310 wall elements} carrying
\textbf{6 layups} (\texttt{layup\_0} through \texttt{layup\_5}) built from
\textbf{6 materials} (\texttt{gelcoat}, \texttt{glass\_triax},
\texttt{glass\_uniax}, \texttt{medium\_density\_foam}, \texttt{carbon\_uniax},
\texttt{glass\_biax}). The element count is confirmed by the recovery output:
\texttt{iea\_s10\_elem\_plate\_forces.dat} has two header lines and 310 data
rows, one per ring element. The fourth script,
\texttt{4\_sweep\_51\_stations.py}, repeats the same computation for all 51
stations of the blade.

\section{Stage 1: windIO blade to per-station 1-D shell SG YAMLs}
\label{sec:windio:stage1}

\subsection*{What it does}

\texttt{OpenSG\_io} is the input layer, and it is a \emph{separate repository}:
\url{https://github.com/bagla0/OpenSG_io}, documented at
\url{https://bagla0.github.io/OpenSG_io/}. It reads a windIO blade description
-- auto-detecting windIO v1 and v2 -- interpolates the outer shape, the layup
stacks and the material database to a requested span fraction, and writes the
cross-section in the two forms an MSG analysis can use:

\begin{itemize}
  \item the \textbf{1-D shell SG YAML} that \texttt{opensg\_shell} reads
        directly, which is what stages 2 and 5 of this chapter consume; and
  \item a \textbf{PreVABS XML} that meshes to a 2-D solid cross-section, which
        is the route that produces the independent VABS reference used to
        validate the recovery in stage 5.
\end{itemize}

It also reads PreVABS XML cross-sections and OpenFAST ElastoDyn/BeamDyn blade
data, so published beam properties can be pulled in as a validation reference.
Its own README is explicit that it does not run windIO or OpenFAST; it reads
those formats so you do not have to hand-build the structure gene.

\subsection*{Installing it}

\begin{lstlisting}[style=shellst]
git clone --recurse-submodules https://github.com/bagla0/OpenSG_io
cd OpenSG_io
pip install -r requirements.txt
python scripts/fetch_prevabs.py
\end{lstlisting}

The \texttt{requirements.txt} in that repository lists \texttt{numpy},
\texttt{scipy}, \texttt{pyyaml} and \texttt{meshio} as the core dependencies,
plus \texttt{gmsh}, \texttt{tetgen} and \texttt{pyvista} for the mesh
generators. A windIO \emph{v2} input file additionally needs the
\texttt{windIO} package installed; v1 files do not. The
\texttt{fetch\_prevabs.py} script downloads the PreVABS release binary for your
operating system, which is required only for the 2-D solid route -- the 1-D
shell YAML needed for stages 2 through 5 is written by pure Python.

\subsection*{The command line}

One CLI handles any of the supported inputs. It writes a set of stations in a
single call:

\begin{lstlisting}[style=shellst]
python scripts/opensg_io.py BAR_URC.yaml out --name bar --stations 0.3 0.5 0.7 --solid
\end{lstlisting}

Here \texttt{BAR\_URC.yaml} is the windIO blade, \texttt{out} is the output
directory, \texttt{-{}-name bar} is the file-name stem, \texttt{-{}-stations}
takes the span fractions $r/R$, and \texttt{-{}-solid} additionally emits the
PreVABS XML and runs PreVABS to build the 2-D solid section. Omit
\texttt{-{}-solid} if you only want the 1-D shell YAMLs.

\subsection*{The Python API}

The API is what you script a 51-station blade with, and it is what generated
the station files used in the rest of this chapter:

\begin{lstlisting}[style=opensg,caption={Building one station's 1-D shell SG and its PreVABS twin.}]
from opensg_io import load_blade, build_cross_section, emit_opensg_yaml, emit_prevabs

blade = load_blade("IEA-22-280-RWT.yaml")      # windIO v1 or v2, auto-detected
cs = build_cross_section(blade, r=0.5)         # one span fraction
emit_opensg_yaml(cs, "shell_r050.yaml")        # the 1-D shell SG for opensg_shell
emit_prevabs(cs, "prevabs_r050", name="r050")  # PreVABS XML -> 2-D solid route
\end{lstlisting}

Loop over your station list and you have the \texttt{iea\_s00\_shell.yaml}
through \texttt{iea\_s50\_shell.yaml} directory that stage 2 points at.

\subsection*{What the 1-D shell SG YAML contains}

\texttt{iea\_s10\_shell.yaml} is the reference for the format. Its top-level
keys, in the order they appear in the file, are listed in
Table~\ref{tab:windio:yamlkeys}.

\begin{table}[htbp]
\centering
\small
\caption{Top-level keys of the 1-D shell SG YAML, in file order.}
\label{tab:windio:yamlkeys}
\begin{tabular}{ll}
\toprule
Key & Content \\
\midrule
\texttt{nodes} & contour node coordinates \texttt{[y1, y2, y3]} in metres, one per line \\
\texttt{elements} & two-node line segments along the section contour \\
\texttt{sets} & \texttt{element} sets, one per layup, holding the element labels \\
\texttt{sections} & \texttt{type: shell}, \texttt{elementSet: layup\_k}, and the \texttt{layup} as \\
  & \texttt{[material, thickness (m), angle (deg)]} triples, outermost ply first \\
\texttt{elementOrientations} & one $3\times3$ frame per element, flattened; the first row is $\mathbf{e}_1$, \\
  & out of the section plane, along the beam axis \\
\texttt{materials} & \texttt{density} in kg/m$^3$ and orthotropic \texttt{elastic: \{E, G, nu\}} triples in Pa \\
\texttt{reference} & which surface the section geometry refers to \\
\bottomrule
\end{tabular}
\end{table}

A concrete section block from that file, the spar-cap laminate:

\begin{lstlisting}[style=shellst]
sections:
- type: shell
  elementSet: layup_0
  layup:
  - - gelcoat
    - 0.0005
    - 0.0
  - - glass_triax
    - 0.00300294
    - 0.0
  - - glass_uniax
    - 0.02609467
    - 0.0
  - - glass_triax
    - 0.00300294
    - 0.0
\end{lstlisting}

The \texttt{reference} field is load-bearing and is the single source of truth
for the whole chapter. \texttt{build\_rm\_bundle} reads it when no override is
passed, and it fixes the through-thickness offset \texttt{frac} used by the
ring laminate law, the plate SG's reference plane, the emitted ABD and the
recovery depth conversion alike. The IEA stations all carry
\texttt{reference: center}.

\begin{table}[htbp]
\centering
\small
\caption{The \texttt{reference} field and the through-thickness fraction it sets.}
\label{tab:windio:reference}
\begin{tabular}{lll}
\toprule
\texttt{reference} & \texttt{frac} & Meaning \\
\midrule
\texttt{center} & 0.5 & mid-surface (what the IEA stations use) \\
\texttt{oml} & 0.0 & outer mould line \\
\texttt{oml\_flip} & 1.0 & outer mould line, reversed stacking \\
\texttt{iml} & 1.0 & inner mould line \\
\bottomrule
\end{tabular}
\end{table}

\section{Stage 2: each station to a Timoshenko $6\times6$}
\label{sec:windio:stage2}

\subsection*{What it computes}

The RM shell ring homogenization: a six-degree-of-freedom-per-node shell
element with an element-wise drilling Lagrange multiplier, $\gamma_{23}$ tied
by MITC, and the MSG (Yu 2002 least-squares) wall transverse-shear law
$\mathbf{G}$. The ring warping problem is solved for the four Euler--Bernoulli
modes $\mathbf{V}_0$ and the Timoshenko correction $\mathbf{V}_1$, and the
energy is condensed to the beam $6\times6$ in the VABS strain order
\[
  \boldsymbol{\varepsilon} =
  \begin{bmatrix} \varepsilon_{11} & \gamma_{12} & \gamma_{13} &
                  \kappa_1 & \kappa_2 & \kappa_3 \end{bmatrix}^{\mathsf T},
\]
whose diagonal is $[\,EA\;\;GA_2\;\;GA_3\;\;GJ\;\;EI_2\;\;EI_3\,]$.

\subsection*{What you edit}

One line at the top of \texttt{1\_homo\_beam\_props.py}:

\begin{lstlisting}[style=opensg]
############### User Input #################################
name = "iea_s10"                  # IEA-22 station at r/R = 0.20
############################################################
\end{lstlisting}

\subsection*{The command}

\begin{lstlisting}[style=shellst]
python 1_homo_beam_props.py
\end{lstlisting}

\subsection*{The API, if you drive it yourself}

\begin{lstlisting}[style=opensg,caption={The whole homogenization is one call.}]
import numpy as np
from opensg_shell import build_rm_bundle

B = build_rm_bundle("iea_s10_shell.yaml")
C6 = np.asarray(B["Timo"])        # (6,6) Timoshenko stiffness
print(B["ref"], B["g_source"])    # 'center'  'msg'
\end{lstlisting}

The full signature is
\texttt{build\_rm\_bundle(shell\_yaml, ref=None, shear="mitc4\_g23",
g\_source="msg")}. Passing \texttt{ref=None} means the YAML's own
\texttt{reference} field decides, which is why homogenization and
dehomogenization cannot silently disagree; pass an explicit \texttt{ref} only
to override. \texttt{g\_source} chooses the wall transverse-shear law,
\texttt{"msg"} for the Yu 2002 least-squares projection or \texttt{"whitney"}
for the complementary-energy shear flow.

What comes back is a bundle, not just a matrix. It carries \texttt{Timo}, the
warping modes \texttt{V0} and \texttt{V1}, the ring geometry
(\texttt{corners}, \texttt{red\_cells}, \texttt{rx3}, \texttt{re3},
\texttt{k22}, \texttt{ax}, \texttt{cross}, \texttt{strip}),
\texttt{layup\_per\_elem}, the geometry-free \texttt{layup\_db} and
\texttt{material\_db}, and \texttt{frac}, \texttt{ref}, \texttt{g\_source}.
Stage 5 consumes every one of those entries, so do not throw the bundle away
after reading the matrix out of it.

\subsection*{What comes out}

\begin{table}[htbp]
\centering
\small
\caption{Files written by a single-station homogenization run.}
\label{tab:windio:stage2out}
\begin{tabular}{lll}
\toprule
File & Written by & Content \\
\midrule
\texttt{iea\_s10\_shell\_Timo.out} & \texttt{build\_rm\_bundle} & the timed SwiftComp \texttt{.K} layout: the banner \\
  &  & \texttt{OpenSG msg-shell beam model [ext sh2 sh3} \\
  &  & \texttt{twist bend2 bend3]}, \texttt{The Effective Timoshenko} \\
  &  & \texttt{Stiffness Matrix}, its compliance, and a \\
  &  & \texttt{Time taken} footer \\
\texttt{iea\_s10\_shell\_ABDG.out} & \texttt{build\_rm\_bundle} & one block per section: \texttt{The Effective Reissner-} \\
  &  & \texttt{Mindlin Plate Stiffness Matrix}, the $8\times8$ ABD \\
  &  & plus $\mathbf{G}$, and its compliance \\
\texttt{abd/iea\_s10\_shell\_abd.yaml} & \texttt{build\_rm\_bundle} & the per-station ABD YAML, emitted once and \\
  &  & cached; carries \texttt{mass\_per\_area} per layup \\
\texttt{iea\_s10\_beam\_Timo.out} & the script & the bare $6\times6$ as text, with the strain order \\
  &  & and the reference recorded in the header \\
\texttt{iea\_s10\_ring\_mesh.png} & the script & the real computed ring, elements coloured by layup \\
\bottomrule
\end{tabular}
\end{table}

Two details about the ABD YAML are worth knowing before you reuse it. It is
written by \texttt{emit\_station\_abd} with \texttt{ref="mid"} when the station
is \texttt{center}-referenced, so its \texttt{reference: mid} field means the
same surface as the ring's \texttt{center}. And it is written with the
emitter's default transverse-shear source, so the file records
\texttt{g\_source: whitney}: the $8\times8$ block in that YAML uses the Whitney
shear-flow $\mathbf{G}$, whereas the ring homogenization itself used the MSG
$\mathbf{G}$. The membrane and bending blocks are identical either way; only
the two transverse-shear rows differ. For \texttt{iea\_s10} the file reports
\texttt{n\_layups: 6}, and for \texttt{layup\_0} a total \texttt{thickness} of
$0.03260055$~m and a \texttt{mass\_per\_area} of $62.89$~kg/m$^2$.

\subsection*{What the numbers mean}

Table~\ref{tab:windio:c6} is the content of \texttt{iea\_s10\_beam\_Timo.out},
in SI units. The YAML is in metres and pascals, so $EA$, $GA_2$ and $GA_3$ are
in N and $GJ$, $EI_2$, $EI_3$ in N$\cdot$m$^2$.

\begin{table}[htbp]
\centering
\footnotesize
\caption{The Timoshenko $6\times6$ of IEA-22 station \texttt{iea\_s10} at $r/R = 0.20$,
         in the VABS strain order, from \texttt{iea\_s10\_beam\_Timo.out}.}
\label{tab:windio:c6}
\setlength{\tabcolsep}{4pt}
\begin{tabular}{lrrrrrr}
\toprule
 & $\varepsilon_{11}$ & $\gamma_{12}$ & $\gamma_{13}$ & $\kappa_1$ & $\kappa_2$ & $\kappa_3$ \\
\midrule
$\varepsilon_{11}$ & 2.76606e+10 & $-$1.72307e$-$05 & 4.66505e$-$06 & 5.40202e$-$06 & 1.66864e+09 & $-$2.07730e+09 \\
$\gamma_{12}$ & $-$1.72307e$-$05 & 7.18613e+08 & 5.75204e+07 & $-$1.01920e+08 & 2.30800e$-$06 & 5.27387e$-$06 \\
$\gamma_{13}$ & 4.66505e$-$06 & 5.75204e+07 & 4.21744e+08 & $-$1.62178e+08 & 8.48168e$-$06 & 2.79314e$-$07 \\
$\kappa_1$ & 5.40202e$-$06 & $-$1.01920e+08 & $-$1.62178e+08 & 2.43757e+09 & 1.64496e$-$06 & $-$8.11406e$-$06 \\
$\kappa_2$ & 1.66864e+09 & 2.30800e$-$06 & 8.48168e$-$06 & 1.64496e$-$06 & 3.51439e+10 & $-$5.55712e+09 \\
$\kappa_3$ & $-$2.07730e+09 & 5.27387e$-$06 & 2.79314e$-$07 & $-$8.11406e$-$06 & $-$5.55712e+09 & 6.81297e+10 \\
\bottomrule
\end{tabular}
\end{table}

Read the matrix as physics rather than as six numbers on a diagonal. The
extension--bending entries $1.67\times10^{9}$ and $-2.08\times10^{9}$ say that
the section's tension centre does not sit on the reference axis. The
bending--bending coupling $-5.56\times10^{9}$ couples flap and edge bending, so
the principal bending axes are not the reference axes either. The shear--twist
entries $-1.02\times10^{8}$ and $-1.62\times10^{8}$ locate the shear centre
away from the reference axis as well. The entries printed at $10^{-5}$ against
diagonals of $10^{10}$ are numerically zero: extension does not couple to
transverse shear or to twist in this section, and the solver reproduces that
null result to fifteen orders of magnitude, which is a useful sanity check that
the ring assembled and condensed correctly.

\subsection*{The spanwise table: all 51 stations}

\texttt{4\_sweep\_51\_stations.py} repeats the homogenization for every station
of the blade and, in the same loop, runs the first-order recovery of stage 5.
What you edit is the block at the top:

\begin{lstlisting}[style=opensg]
############### User Input #################################
yaml_dir = os.path.expanduser(
    "~/OpenSG-TW-claude/examples/data/iea_all_stations/shell51/1d_yaml")
ff_file = "ff51_rmc_reform.dat"
n_depth = 9                  # through-thickness recovery points per element
stations = range(51)         # s00 .. s50
############################################################
\end{lstlisting}

\texttt{yaml\_dir} is where \texttt{iea\_s00\_shell.yaml} through
\texttt{iea\_s50\_shell.yaml} live, which is the OpenSG\_io output of stage 1.
Point it at your own directory. The script builds each path as
\texttt{iea\_s\%02d\_shell.yaml} against that directory, and reads the load
table \texttt{ff\_file} from the current directory.

\begin{lstlisting}[style=shellst]
python 4_sweep_51_stations.py
\end{lstlisting}

A station whose YAML is missing prints
\texttt{s\%02d: yaml missing -{}- skipped} and is passed over, and any station
that raises is caught and reported as \texttt{s\%02d: FAILED} so that one bad
section cannot kill a sweep that has already spent minutes on the others. In
the shipped run, \textbf{49 of the 51 stations completed}: both
\texttt{iea51\_beam\_props.dat} and \texttt{iea51\_peak\_stress.dat} carry
$\eta = 0.00$ through $0.94$ in steps of $0.02$ plus $\eta = 1.00$, with
$\eta = 0.96$ and $\eta = 0.98$ absent.

One indexing rule matters here. The sweep takes station $i$'s load vector as
row $i$ of the FF table, \texttt{FF = ffall[i, 1:]}, and its span fraction as
\texttt{eta = ffall[i, 0]}. The row order of the load table must therefore match
the station numbering of the YAML files. In the shipped data, station
\texttt{s10} is row 10 and $\eta = 0.20$.

\texttt{iea51\_beam\_props.dat} holds one row per completed station: $\eta$
followed by the six diagonal terms, written with the header
\texttt{IEA-22 spanwise Timoshenko beam properties (RM shell ring)}. Selected
rows, verbatim from the file:

\begin{table}[htbp]
\centering
\footnotesize
\caption{Selected rows of \texttt{iea51\_beam\_props.dat}. Stiffnesses in N
         ($EA$, $GA_2$, $GA_3$) and N$\cdot$m$^2$ ($GJ$, $EI_2$, $EI_3$).}
\label{tab:windio:props}
\setlength{\tabcolsep}{4.5pt}
\begin{tabular}{lrrrrrr}
\toprule
$\eta = r/R$ & $EA$ & $GA_2$ & $GA_3$ & $GJ$ & $EI_2$ & $EI_3$ \\
\midrule
0.00 & 6.44275512e+10 & 7.96306621e+09 & 8.36135930e+09 & 1.27413308e+11 & 2.97295270e+11 & 2.19213722e+11 \\
0.10 & 2.89212767e+10 & 1.11237808e+09 & 1.04545606e+09 & 1.14245218e+10 & 8.40070906e+10 & 1.21502270e+11 \\
0.20 & 2.76605989e+10 & 7.18613262e+08 & 4.21744284e+08 & 2.43756534e+09 & 3.51438873e+10 & 6.81297238e+10 \\
0.50 & 2.15596413e+10 & 5.34263266e+08 & 2.02128533e+08 & 4.95011412e+08 & 7.26524060e+09 & 1.82791086e+10 \\
0.80 & 9.28038409e+09 & 3.90243799e+08 & 7.82870454e+07 & 6.27349440e+07 & 5.76401363e+08 & 2.07536066e+09 \\
0.94 & 2.64020156e+09 & 1.80834082e+08 & 3.43479097e+07 & 1.35795150e+07 & 6.21667876e+07 & 4.27156667e+08 \\
1.00 & 2.65581077e+08 & 1.56743220e+07 & 1.88135914e+06 & 5.15174420e+04 & 3.04364786e+05 & 1.60428557e+06 \\
\bottomrule
\end{tabular}
\end{table}

The $\eta = 0.20$ row is digit-for-digit the diagonal of the standalone
\texttt{iea\_s10} run of Table~\ref{tab:windio:c6}: $EA$ is
\texttt{2.76605989e+10} in both, $GJ$ is \texttt{2.43756534e+09} in both, and so
on for all six. The sweep and the single-station script are the same
computation, which is the cheapest available check that your sweep is wired
correctly -- run \texttt{1\_homo\_beam\_props.py} on one station and match its
row.

\begin{figure}[htbp]
\centering
\includegraphics[width=0.95\textwidth]{\FIGDIR/iea51_beam_props.png}
\caption{Spanwise Timoshenko diagonal of the IEA-22 blade: the six terms
         $EA$, $GA_2$, $GA_3$, $GJ$, $EI_2$, $EI_3$ on log axes against $r/R$,
         one panel each, produced by \texttt{4\_sweep\_51\_stations.py}.}
\label{fig:windio:beamprops}
\end{figure}

Figure~\ref{fig:windio:beamprops} plots those six columns. Three regimes are
visible. The root collapses fast as the circular root transitions into an
aerofoil: $GJ$ falls from $1.274\times10^{11}$ at $\eta = 0$ to
$1.142\times10^{10}$ N$\cdot$m$^2$ at $\eta = 0.10$, an order of magnitude in
the first tenth of span, and $GA_2$ falls from $7.96\times10^{9}$ to
$1.11\times10^{9}$ N over the same interval. The mid-span then decays smoothly
as the section thins. The tip drops off a cliff. Across the whole blade $GJ$
goes from $1.27\times10^{11}$ to $5.15\times10^{4}$ N$\cdot$m$^2$, six and a
half decades, which is why the figure is drawn on log axes and why a beam
solver must be given the table rather than an average.

\section{Stage 3: handing the table to BeamDyn}
\label{sec:windio:stage3}

Be explicit about what changes hands here. \textbf{OpenSG produces the sectional
$6\times6$ stiffness matrix, and nothing else.} BeamDyn, the geometrically exact
beam module of OpenFAST, is a separate program that you install, configure and
run yourself. OpenSG-2.0 ships no BeamDyn input file and no BeamDyn results, and
this chapter does not attempt to document BeamDyn's input format: for the
meaning, ordering and units of every field in a BeamDyn blade input file, and
for the choice of quadrature and the number of input stations, consult the
OpenFAST BeamDyn documentation, which is the authority on that file.

\subsection*{The bridge that OpenSG\_io provides}

\texttt{OpenSG\_io} contains an \texttt{openfast\_io} module whose job is
exactly this hand-off. Two functions write the OpenFAST side:

\begin{lstlisting}[style=opensg,caption={Writing a BeamDyn blade input from a spanwise table of OpenSG 6x6 matrices.}]
import numpy as np
from opensg_io.openfast_io import timo_to_beamdyn, write_beamdyn_blade

# etas: (n,) span fractions; Ks: list of n OpenSG Timoshenko 6x6 matrices
Kbd = timo_to_beamdyn(Ks[0])                 # one matrix, reordered
write_beamdyn_blade(etas, Ks, M_list=None,
                    out_path="blade_BeamDyn_Blade.dat")
\end{lstlisting}

\texttt{timo\_to\_beamdyn(Ktimo)} reorders an OpenSG Timoshenko $6\times6$,
which is stored in the VABS order $[EA, GA_2, GA_3, GJ, EI_2, EI_3]$, into the
BeamDyn degree-of-freedom order. It applies a symmetric permutation
\texttt{Ktimo[np.ix\_(q, q)]} with \texttt{q = [1, 2, 0, 4, 5, 3]}, which is the
mapping in Table~\ref{tab:windio:bdorder}. Because the permutation is applied to
both rows and columns, every off-diagonal coupling is carried across intact.

\begin{table}[htbp]
\centering
\small
\caption{The reordering performed by \texttt{timo\_to\_beamdyn}: which OpenSG
         diagonal term lands in which BeamDyn slot.}
\label{tab:windio:bdorder}
\begin{tabular}{clc}
\toprule
BeamDyn slot & Meaning & From OpenSG term \\
\midrule
1 & shear\_x & $GA_2$ \\
2 & shear\_y & $GA_3$ \\
3 & axial & $EA$ \\
4 & bend\_x & $EI_2$ \\
5 & bend\_y & $EI_3$ \\
6 & torsion & $GJ$ \\
\bottomrule
\end{tabular}
\end{table}

\texttt{write\_beamdyn\_blade(etas, K\_timo\_list, M\_list=None,
out\_path="blade\_BeamDyn\_Blade.dat", damp=(0.003, 0.002, 0.002, 0.002, 0.003,
0.002))} takes the span fractions and the list of OpenSG matrices, reorders each
one internally with \texttt{timo\_to\_beamdyn}, and writes a BeamDyn blade input
with a \texttt{station\_total} equal to the number of entries you supplied,
\texttt{damp\_type} set to $1$, the six damping coefficients from
\texttt{damp}, and then a stiffness block followed by a mass block for each
station. Verify the resulting file against the BeamDyn documentation for the
OpenFAST version you are running before trusting it.

\subsection*{The mass matrix is yours}

\texttt{write\_beamdyn\_blade} takes \texttt{M\_list} as an \emph{optional}
argument and \textbf{writes a $6\times6$ block of zeros for every station when it
is omitted}. That is a placeholder, not a mass model. Nothing in this pipeline
computes a sectional mass matrix: the only mass quantity OpenSG-2.0 produces
here is \texttt{mass\_per\_area} per laminate in the emitted ABD YAML
(Section~\ref{sec:windio:stage2}). Building the sectional inertia matrix,
assembling the complete BeamDyn input, and running the aeroelastic simulation
are all yours, outside OpenSG.

\subsection*{Use the full matrix, not the table's diagonal}

\texttt{iea51\_beam\_props.dat} records only the six diagonal terms, for
plotting and inspection. The couplings displayed in Table~\ref{tab:windio:c6}
are precisely what a composite blade beam model exists to carry, so feed the
beam solver the full $6\times6$ of each station -- either from that station's
\texttt{<station>\_Timo.out} or, more conveniently, straight from
\texttt{B["Timo"]} in your own driver script, collected into the
\texttt{K\_timo\_list} above.

\subsection*{Reading OpenFAST data back as a reference}

The same module reads in the other direction, which is how you validate an
OpenSG table against a published blade. \texttt{read\_beamdyn\_blade(path)}
parses a BeamDyn blade input and returns \texttt{(etas, K\_list, M\_list)} with
each matrix in BeamDyn order; \texttt{beamdyn\_to\_timo(Kbd)} reorders one back
into the OpenSG order so it can be compared term by term.
\texttt{read\_elastodyn\_blade(path)} parses an ElastoDyn blade file and returns
arrays of \texttt{BlFract}, \texttt{PitchAxis}, \texttt{StrcTwst} in degrees,
\texttt{BMassDen} in kg/m, and \texttt{FlpStff} and \texttt{EdgStff} in
N$\cdot$m$^2$, which map to OpenSG's $EI_2$ (flap) and $EI_3$ (edge);
\texttt{elastodyn\_at(ed, r)} interpolates those to any span fraction.

\section{Stage 4: the loads come back as the FF table}
\label{sec:windio:stage4}

A BeamDyn or GEBT run returns, at every station, the sectional force and moment
resultants the blade actually carries in a given load case. That table is the
input to recovery. \texttt{ff51\_rmc\_reform.dat} is such a table. It has one
comment header line and 51 data rows, one per station. The header and the first
data rows, verbatim:

\begin{lstlisting}[style=shellst]
# eta F1 F2 F3 M1 M2 M3 (RM CENTER-ref BeamDyn, VABS order)
0.00000000e+00 -6.44900592e-06 -2.39422988e-08 2.23607775e+06 1.21251850e+06 -1.12030544e+08 -4.11329552e-07
2.00000000e-02 2.07495361e+03 7.56198425e+01 2.15497000e+06 1.21427150e+06 -1.05899552e+08 2.54692651e+03
4.00000000e-02 4.74173926e+03 1.68900665e+02 2.07209938e+06 1.23365012e+06 -9.99422000e+07 5.32342041e+03
\end{lstlisting}

Column 1 is $\eta = r/R$, running from $0.00$ to $1.00$ in steps of $0.02$.
Columns 2 through 4 are the forces $F_1$, $F_2$, $F_3$ in N; columns 5 through 7
are the moments $M_1$, $M_2$, $M_3$ in N$\cdot$m. The row used by the
single-station scripts is row 10:

\begin{lstlisting}[style=shellst]
2.00000000e-01 4.19268867e+04 8.26739746e+03 1.46733512e+06 1.21914500e+06 -6.01849720e+07 3.04264594e+05
\end{lstlisting}

The frame and the component order stated in that header are not decoration. The
resultants must be expressed about the same reference axis and in the same
component order as the $6\times6$ that produced them, or the round trip does not
close and the recovered stress is silently wrong. Here both are \emph{RM
center-ref} and VABS order, matching \texttt{reference: center} in the station
YAMLs. Note also the $\eta = 0$ row: its axial and in-plane entries print as
$-6.4\times10^{-6}$ and $-2.4\times10^{-8}$ N, a numerically zero root row
rather than a physical load.

A converter turns the table into an annotated YAML and pulls one station out of
it:

\begin{lstlisting}[style=opensg,caption={ff51\_rmc\_reform.yaml: the FF table as a self-describing file.}]
from opensg_shell.helper.ff_to_yaml import convert, station_ff

d = convert("ff51_rmc_reform.dat")                      # writes ff51_rmc_reform.yaml
FF, eta = station_ff("ff51_rmc_reform.yaml", eta=0.20)  # (6,) resultants, and 0.2
\end{lstlisting}

\texttt{convert(dat\_path, out\_base=None, frame="RM center-ref")} reads the
seven-column table, raising if it has fewer than seven columns, and writes
\texttt{<base>.yaml}. \texttt{station\_ff(yaml\_path, eta=None, row=None)}
returns the $(6,)$ resultant vector and the $\eta$ of one station, selected
either by nearest $\eta$ or by row index. The resulting
\texttt{ff51\_rmc\_reform.yaml} records the convention explicitly next to the
data, which is what makes the file readable months later:

\begin{lstlisting}[style=shellst]
source: ff51_rmc_reform.dat
convention:
  order: [F1, F2, F3, M1, M2, M3]
  frame: RM center-ref
  units: {force: N, moment: N.m}
stations:
- row: 0
  eta: 0.0
  FF: [-6.44900592e-06, -2.39422988e-08, 2236077.75, 1212518.5, -112030544.0, -4.11329552e-07]
- row: 10
  eta: 0.2
  FF: [41926.8867, 8267.39746, 1467335.12, 1219145.0, -60184972.0, 304264.594]
\end{lstlisting}

\section{Stage 5: dehomogenization back to 3-D wall stress}
\label{sec:windio:stage5}

\subsection*{What it computes}

A two-step chain, because a blade wall is a laminate inside a section:
\[
  \text{beam } \mathbf{FF}
  \;\xrightarrow[\text{RM shell ring}]{}\;
  \text{per-element plate forces}
  \;\xrightarrow[\text{RM plate SG}]{}\;
  \boldsymbol{\sigma}(z)\ \text{through the wall.}
\]

\emph{Step 5a, the section} (\texttt{opensg\_shell}). The macro beam strain is
$\bar{\boldsymbol{\varepsilon}} = \mathbf{C}_6^{-1}\mathbf{FF}$;
\texttt{\_macro\_fields(B, beam\_force\_vabs=FF)} solves it and recombines the
nodal RM warping fields $\mathbf{w} = \mathbf{V}_0\bar{\boldsymbol{\varepsilon}}
+ \mathbf{V}_1\boldsymbol{\varepsilon}'$ and their derivatives, and
\texttt{\_rm\_shell\_strain(B, e, xi, st\_m, aA, aB)} evaluates, for ring
element \texttt{e} at arc coordinate \texttt{xi}, the six shell strains
\[
  \mathbf{s}_6 = \begin{bmatrix} \varepsilon_{11} & \varepsilon_{22} &
  2\varepsilon_{12} & \kappa_{11} & \kappa_{22} & 2\kappa_{12}\end{bmatrix}.
\]
That wall's own MSG ABD turns them into the \textbf{plate reaction forces}
\[
  \mathbf{F}_6 = \mathbf{A}_6\,\mathbf{s}_6 =
  \begin{bmatrix} N_{11} & N_{22} & N_{12} & M_{11} & M_{22} & M_{12}\end{bmatrix},
\]
the interface quantity handed down to the plate SG.

\emph{Step 5b, the wall} (\texttt{opensg\_solid.rm\_plate\_1D}).
\texttt{rm\_plate\_msg(thick, angles\_deg, mat\_names, material\_db,
fraction=frac)} builds the through-thickness MSG-RM structure gene for each
layup once, and \texttt{msgrm\_strain\_at\_depth} -- or its vectorized twin
\texttt{msgrm\_strain\_at\_depth\_batch}, which evaluates many points of one
laminate in a handful of einsums -- returns the 3-D Voigt strain and stress at
depth $z$. Both default to \texttt{frame="material"}, meaning the result is
rotated into the ply material frame where failure criteria live; pass
\texttt{frame="plate"} to stay in the wall frame.

Recovery needs strain \emph{gradients}, and they come from the ring itself. The
span gradient $\partial\mathbf{E}/\partial x_1$ comes from the macro recovery
ladder, computed per element as \texttt{BDe @ st\_cl1 + BDh @ aB[g]} with the
element operators from \texttt{quad\_ops\_indep}. The arc gradient
$\partial\mathbf{E}/\partial x_2$ comes from nodal differencing of
$\mathbf{s}_6$ between the two elements sharing a node -- but only where those
two elements carry the \emph{same} layup, because a layup change is a physical
jump and must never be differenced across. Passing the second gradients as well
(\texttt{dE11}, \texttt{dE12}, \texttt{dE22}) selects the second-order (V2)
recovery, the rung that can in principle carry the transverse pair
$\sigma_{33}$, $\sigma_{13}$ that first order cannot.

\subsection*{What you edit}

The user block at the top of \texttt{2\_dehom\_v2\_vabs.py}:

\begin{lstlisting}[style=opensg]
############### User Input #################################
name = "iea_s10"                     # IEA-22 station at r/R = 0.20
ff_file = "ff51_rmc_reform.dat"      # BeamDyn FF (RM center-ref, VABS order)
station_row = 10                     # row 10 -> eta = 0.20
vabs_sm = "iea_s10.sg.SM"            # the VABS 2-D solid reference cloud
jpath_file = "iea_s10.lp_sparcap_left_thickness.coords"
############################################################
\end{lstlisting}

\begin{lstlisting}[style=shellst]
python 2_dehom_v2_vabs.py
\end{lstlisting}

\subsection*{What comes out}

\begin{table}[htbp]
\centering
\small
\caption{Files written by \texttt{2\_dehom\_v2\_vabs.py}.}
\label{tab:windio:stage5out}
\begin{tabular}{ll}
\toprule
File & Content \\
\midrule
\texttt{iea\_s10\_report.txt} & the validation report: the FF vector, the Timo diagonal, rms \\
  & differences against VABS split web/skin, the V2 payload block, \\
  & and the through-thickness path tables \\
\texttt{iea\_s10\_elem\_plate\_forces.dat} & 310 rows, one per ring element: element index, \texttt{y2 y3}, the six \\
  & shell strains, and the six plate reaction forces \\
  & $N_{11}, N_{22}, N_{12}$ [N/m] and $M_{11}, M_{22}, M_{12}$ [N] \\
\texttt{iea\_s10\_dehom\_v2.txt} & 166\,302 rows, one per VABS gauss point: \texttt{y2 y3 elem z}, then \\
  & the six recovered V2 stresses and the six VABS stresses, all MPa \\
\texttt{iea\_s10\_dehom\_v2.npz} & the same fields plus \texttt{C6}, \texttt{FF}, \texttt{s6mid}, \texttt{F6mid} and every strain \\
  & gradient, for re-plotting without re-solving \\
\texttt{iea\_s10\_v2\_section.png} & section contours, VABS beside RM V2, for $\sigma_{11}$, $\sigma_{12}$, $\sigma_{13}$, $\sigma_{33}$ \\
\texttt{iea\_s10\_parity.png} & parity clouds, first order and V2 against VABS \\
\texttt{iea\_s10\_thickness\_paths.png} & through-thickness profiles on a spar-cap wall, a web, and the \\
  & station's own reference path \\
\bottomrule
\end{tabular}
\end{table}

The first two data rows of \texttt{iea\_s10\_elem\_plate\_forces.dat} show what
the interface quantity looks like in practice:

\begin{lstlisting}[style=shellst]
# per-element PLATE REACTION FORCES from the RM shell beam model (Timo), eta=0.20
# elem  y2(m) y3(m)  e11 e22 2e12 k11 k22 2k12  N11 N22 N12(N/m) M11 M22 M12(N)
     0  4.677080e+00  1.048036e-01  5.426575e-04 -1.683049e-04  5.685206e-04 ...
\end{lstlisting}

\subsection*{What it was validated against}

\texttt{iea\_s10.sg.SM} is the VABS stress cloud of the \emph{same}
cross-section meshed as a 2-D solid, which is the independent reference: the
recovery is projected onto its gauss points and differenced there, 166\,302
points in all. Every point is assigned to its nearest ring element, a local arc
coordinate and a signed wall depth, and the VABS columns are permuted into the
$(11, 22, 33, 23, 13, 12)$ order and converted from Pa to MPa before the
comparison. Table~\ref{tab:windio:vabs} is copied from
\texttt{iea\_s10\_report.txt}, with the VABS rms alongside so each difference
has a scale.

\begin{table}[htbp]
\centering
\small
\caption{rms difference against VABS in MPa, web / skin, from
         \texttt{iea\_s10\_report.txt} at $\eta = 0.20$.}
\label{tab:windio:vabs}
\begin{tabular}{lrrr}
\toprule
Component & first order (web / skin) & V2 (web / skin) & VABS rms (web / skin) \\
\midrule
$\sigma_{11}$ & 43.073 / 9.529 & 43.073 / 9.529 & 44.440 / 96.676 \\
$\sigma_{22}$ & 0.612 / 0.482 & 0.612 / 0.482 & 0.704 / 0.998 \\
$\sigma_{12}$ & 6.799 / 1.255 & 6.798 / 1.255 & 22.026 / 6.009 \\
$\sigma_{33}$ & 0.081 / 0.049 & 0.081 / 0.049 & 0.081 / 0.049 \\
$\sigma_{13}$ & 0.394 / 0.192 & 0.394 / 0.192 & 0.394 / 0.193 \\
$\sigma_{23}$ & 0.035 / 0.028 & 0.035 / 0.028 & 0.035 / 0.028 \\
\bottomrule
\end{tabular}
\end{table}

Read this honestly, because the table says three different things at once.

\emph{The in-plane skin stress is good.} $\sigma_{11}$ in the skin differs by
9.53 MPa rms against a VABS rms of 96.68 MPa, and $\sigma_{12}$ by 1.26 against
6.01. The report also tabulates three through-thickness paths, and those
sharpen the picture considerably. On the station's own reference thickness path,
whose 15 points sit \textbf{0.00 mm} from the nearest VABS gauss point and are
therefore exactly in the reference frame, the $\sigma_{11}$ difference is
\textbf{1.265 MPa against a VABS rms of 286.837 MPa}, better than half a
percent. The two paths generated from the ring geometry sit 2.44 mm and 2.43 mm
from the nearest VABS point, and their larger differences -- 38.355 MPa rms on
the spar cap, 2.299 MPa on the web -- are dominated by that geometric offset
rather than by the recovery itself. The printed nearest-VABS distance is the
frame check: millimetres means the same frame, hundreds of millimetres means you
are comparing against a foreign mesh.

\emph{The web is the weak spot.} For $\sigma_{11}$ in the web the difference,
43.07, is essentially the VABS rms itself, 44.44, which means the web axial
stress is not being reproduced. The same pattern holds for the three transverse
components everywhere: the difference equals the VABS rms to three digits, which
is what you get when the recovered field is very nearly zero. This is the known
transverse recovery gap of the shell-to-plate route, not a tuning problem.

\emph{The second-order rung is inert at this load state.} The report's V2
payload block measures what V2 \emph{changes}, rather than a rounded difference
that could hide a small effect. At this station $\sigma_{11}$ moves by
$1.950\times10^{-4}$ MPa rms against a first-order rms of $8.851\times10^{1}$,
$\sigma_{12}$ by $5.685\times10^{-3}$ against $1.068\times10^{1}$,
$\sigma_{33}$ by $5.190\times10^{-5}$ MPa against a VABS rms of
$5.576\times10^{-2}$, and $\sigma_{13}$ and $\sigma_{23}$ by exactly
$0.000\times10^{0}$. For this constant-force blade state the V2 drivers carry no
payload, and the first-order $\sigma_{33}$ rms is $3.603\times10^{-13}$ MPa,
machine zero.

\subsection*{The correctness check on transverse shear}

\texttt{3\_q\_consistency\_check.py} tests an identity that any recovered
transverse shear must satisfy, wall by wall: the thickness integral of the
recovered stress must equal the wall's own transverse-shear resultant, which the
RM shell already knows.
\[
  I_a = \int_h \sigma_{a3}\,\mathrm{d}z
  \;\;\overset{!}{=}\;\;
  Q_a = \left(\mathbf{G}_{\mathrm{msg}}\,\mathbf{s}_2\right)_a ,
  \qquad
  \mathbf{s}_2 = \begin{bmatrix} 2\gamma_{13} & 2\gamma_{23}\end{bmatrix}.
\]
The plate pipeline enforces this by construction -- shape from the strain
gradients, amplitude from the resultant, $\sigma_{a3}$ rescaled by $Q/I$. The
shell-to-plate chain never uses $\mathbf{s}_2$, so nothing pins the amplitude,
and the ratio $r = I/Q$ measures how far off it is. Consistency means $r = 1$.

\begin{lstlisting}[style=opensg]
############### User Input #################################
name = "iea_s10"
ff_file = "ff51_rmc_reform.dat"
station_row = 10                  # eta = r/R = 0.20
n_z = 41                          # integration points through each wall
############################################################
\end{lstlisting}

\begin{lstlisting}[style=shellst]
python 3_q_consistency_check.py
\end{lstlisting}

The script writes \texttt{iea\_s10\_q\_consistency.txt} and
\texttt{iea\_s10\_q\_consistency.png}, the latter plotting recovered $|I_{13}|$
against wall $|Q_1|$ on log axes with the consistency line, beside the histogram
of $\log_{10}(I_{13}/Q_1)$. The measured result:

\begin{table}[htbp]
\centering
\small
\caption{Q-consistency ratios at $\eta = 0.20$, from
         \texttt{iea\_s10\_q\_consistency.txt}.}
\label{tab:windio:qcheck}
\begin{tabular}{lrrrrl}
\toprule
Ratio & median & mean & min & max & Verdict \\
\midrule
$I_{13}/Q_1$ & 5.228e$-$01 & $-$6.974e+00 & $-$1.299e+03 & 2.260e+01 & target invalid \\
$I_{23}/Q_2$ & 9.895e$-$01 & 1.096e+00 & $-$1.217e+01 & 1.194e+01 & target valid \\
\bottomrule
\end{tabular}
\end{table}

\textbf{$\sigma_{23}$ passes.} The median $I_{23}/Q_2$ is $0.9895$, a
\textbf{1.05\% Q-consistency error}, with a wall $Q_2$ rms of $8.9686$ N/m
against a recovered $I_{23}$ rms of $8.3085$ N/m.

\textbf{$\sigma_{13}$ is not measurable in a cross-section model}, and the
script explains why rather than papering over it. This ring is homogenized with
\texttt{shear="mitc4\_g23"}: $\gamma_{23}$ is tied by MITC while $\gamma_{13}$
is raw and carries the flat-wall drilling mode, so a $Q_1$ built from it is not
a physical wall shear and cannot serve as a recovery target. The fingerprint is
in the magnitudes -- $Q_1$ rms is $9.9064\times10^{2}$ N/m against a $Q_2$ rms
of $8.9686$ N/m, a factor of about 110, which is implausible for a thin-walled
section that carries its load as in-plane shear flow. The ring is a
one-quad-deep prismatic strip, so the $\gamma_{13}$ tying interpolates along a
direction in which the row is already linear and reproduces it exactly: the
script measures \texttt{rms |Q1(g23) - Q1(both)| = 0.0000e+00 N/m} and the same
for $Q_2$, so \texttt{mitc4\_g23} and \texttt{mitc4\_both} are bit-identical
here even though both flags are plumbed and branch. A physical $\gamma_{13}$
needs the surface-quad route, where both rows are genuinely tied, not a rescale
of this ring against a contaminated $Q_1$.

The practical rule that follows: trust $\sigma_{11}$, $\sigma_{22}$ and
$\sigma_{12}$ from this pipeline, treat $\sigma_{23}$ as good to about one
percent in amplitude, and do not use the cross-section route's $\sigma_{13}$ or
$\sigma_{33}$ as a design quantity.

\subsection*{The same recovery, all 51 stations}

\texttt{4\_sweep\_51\_stations.py} runs the \emph{first-order} recovery -- it
passes only the two first gradients to
\texttt{msgrm\_strain\_at\_depth\_batch} -- at \texttt{n\_depth = 9} depths per
wall element. The depths are mid-bin in the normalized thickness coordinate,
$\zeta = (k + \tfrac12)/9$ for $k = 0 \ldots 8$, converted to a physical depth
by $z = (\zeta - \texttt{frac})\,h$ with $h$ the total thickness of that
element's layup and \texttt{frac} taken from the bundle. The peak $|\sigma|$ per
component over the whole section, in the ply material frame, goes to
\texttt{iea51\_peak\_stress.dat}.

\begin{table}[htbp]
\centering
\footnotesize
\caption{Selected rows of \texttt{iea51\_peak\_stress.dat}: peak $|\sigma|$ per
         component over the whole section, ply frame, in MPa. Column order
         follows the file header.}
\label{tab:windio:peak}
\setlength{\tabcolsep}{4.5pt}
\begin{tabular}{lrrrrrr}
\toprule
$\eta = r/R$ & $\sigma_{11}$ & $\sigma_{22}$ & $\sigma_{33}$ & $\sigma_{23}$ & $\sigma_{13}$ & $\sigma_{12}$ \\
\midrule
0.00 & 1.49271417e+02 & 1.53636678e+00 & 4.13972884e$-$13 & 5.39458803e$-$03 & 4.12768841e$-$02 & 2.24186054e+00 \\
0.16 & 3.25793641e+02 & 4.14543750e+00 & 8.97794962e$-$13 & 3.91335971e$-$02 & 3.51941097e$-$02 & 7.01052831e+00 \\
0.20 & 3.18177433e+02 & 3.08795353e+00 & 9.12696123e$-$13 & 1.57149557e$-$02 & 3.45770096e$-$02 & 5.66031217e+00 \\
0.54 & 2.14382707e+02 & 1.62953363e+00 & 8.13510269e$-$13 & 1.74810053e$-$02 & 8.23774625e$-$02 & 5.64234617e+01 \\
0.80 & 1.11316749e+02 & 1.94984342e+00 & 1.93947926e$-$13 & 1.15167225e$-$02 & 5.58043894e$-$02 & 3.85749086e+01 \\
1.00 & 4.88165692e$-$01 & 2.71026227e$-$02 & 2.03726813e$-$15 & 4.70632777e$-$05 & 5.13380478e$-$02 & 7.89278421e+00 \\
\bottomrule
\end{tabular}
\end{table}

Reading the whole file rather than these six rows: the peak axial stress of the
entire blade in this load case is \textbf{325.79 MPa at $\eta = 0.16$}, and the
peak in-plane shear is \textbf{56.42 MPa at $\eta = 0.54$}, where the shear webs
take over from the caps. The $\sigma_{33}$ column never leaves the range
$2\times10^{-15}$ to $2\times10^{-12}$ MPa over all 49 stations -- identically
zero to round-off, the first-order limit noted above -- so do not read it as a
physical through-thickness stress. The $\sigma_{13}$ column, which stays between
$0.033$ and $0.105$ MPa over the whole blade, is subject to the same caveat as in the Q-consistency check and is
not a design quantity either.

\begin{figure}[htbp]
\centering
\includegraphics[width=0.95\textwidth]{\FIGDIR/iea51_sweep.png}
\caption{The two halves of the 51-station sweep. Left: the spanwise peaks of
         $\sigma_{11}$ and $\sigma_{12}$ against $r/R$. Right: the recorded
         per-station wall time, homogenization against dehomogenization.}
\label{fig:windio:sweep}
\end{figure}

Figure~\ref{fig:windio:sweep} shows both halves. The left panel plots the
spanwise peaks of $\sigma_{11}$ and $\sigma_{12}$: the axial peak climbs steeply
out of the root, crests near $\eta \approx 0.16$ and decays to almost nothing at
the tip, while the shear peak rises in distinct steps rather than smoothly, near
$\eta \approx 0.36$ and $\eta \approx 0.46$, where the peak jumps from 4.85 to
25.75 MPa and from 24.40 to 54.30 MPa respectively as a different wall takes
over the shear. The right panel records the per-station wall time of the two
halves separately, so you can see where a sweep's cost actually sits before
scaling it to more stations or more blades. The same arrays are saved to
\texttt{iea51\_sweep.npz} under the keys \texttt{eta}, \texttt{props},
\texttt{peak}, \texttt{t\_homo}, \texttt{t\_dehom} and \texttt{stations}, so a
sweep never has to be repeated just to redraw a figure.

\section{Where each tool ends}
\label{sec:windio:whodoeswhat}

\begin{table}[htbp]
\centering
\small
\caption{Who does what across the three tools in the round trip.}
\label{tab:windio:whodoeswhat}
\begin{tabular}{p{0.30\textwidth}p{0.34\textwidth}p{0.28\textwidth}}
\toprule
OpenSG\_io & OpenSG-2.0 & OpenFAST / BeamDyn \\
\midrule
Reads windIO v1 and v2, PreVABS XML, and OpenFAST ElastoDyn/BeamDyn blade data.
&
Homogenizes each 1-D shell SG to the Timoshenko $6\times6$ and writes
\texttt{<station>\_Timo.out}.
&
Runs the blade dynamics. Installed, configured and run by you, outside OpenSG.
\\[4pt]
Writes the per-station 1-D shell SG YAMLs that \texttt{opensg\_shell} consumes.
&
Writes the RM wall law per section, \texttt{<station>\_ABDG.out}, and the ABD
YAML with \texttt{mass\_per\_area}.
&
Returns the per-station sectional force and moment resultants that come back as
the FF table.
\\[4pt]
Writes the PreVABS XML for the 2-D solid route, which supplies the independent
VABS reference.
&
Recovers the 3-D wall stress from the FF table and validates it against that
VABS cloud.
&
Owns the blade input file format; consult the BeamDyn documentation for it.
\\[4pt]
Bridges the hand-off with \texttt{timo\_to\_beamdyn} and
\texttt{write\_beamdyn\_blade}, and reads OpenFAST data back as a validation
reference.
&
Runs the Q-consistency diagnostic on the recovered transverse shear.
&
Sectional mass matrices are not produced by either OpenSG tool and remain your
responsibility.
\\
\bottomrule
\end{tabular}
\end{table}

The chain closes only because the conventions match at every hand-off: the same
reference surface (\texttt{center}), the same component order (VABS), and the
same units (SI, metres and pascals in, newtons and newton-metres back). Those
three lines are the first things to check when a recovered stress field looks
wrong, and they are checkable without rerunning anything -- the station YAML
records its \texttt{reference}, the FF table records its frame and order in the
header line, and the ABD YAML records the reference it was emitted at.
